% Pursuing Stacks (cotes 134-2 to 134-8): archive cote 134-5, site page 60 (PDF page 61), typescript page 247.
% typescript. Draft compared: ps4.tex:3018-3090.
% AI-transcribed by Claude Opus 5.5 (claude-opus-5-5, adaptive thinking, effort high), Message Batches API, batch ps2 (msgbatch_01NVwr1Yuo8qJf5rykeGBXPK), 2026-11-12; not yet checked by eye.
\scanpage{134-5/60}
\mspage{247}
20.6.

79. I \unsure{would} like now to write out the relationship between asphericity
structures, and contractibility structures as defined in section 51 D)
(pages 117-119). First we'll need to rid ourselves of the smallness as\xout{\unsure{s}}%
sumption for a generating category of an asphericity structure:

\uline{Proposition}\hand{\unsure{1}} Let M be a $\underline{U}$-category ($\underline{U}$ being our basic universe), and N
any full subcategory. The following conditions are e\over{qm}{qu}ivalent:

(i) There exists an asphericity structure $M_{as}$ in M such that
(a) $M_{as}\supset N$ and (b) N admits a \uline{small} full subcategory\inserted{$N_o$} generating
the asphericity structure, i.e. such that the inclusion functor $i_o$ from
$N_o$ to M be $M_{as}$-aspheric.

(ii) There exists a small full subcategory $N_o$ of N, such that for
any x in N, $i_o^{*}(x)$ in $N_o^{\wedge}$ be aspheric, where $i_o : N_o \to M$ is the inclusion
functor.

(iii) The couple (N,N) is an asphericity structure.

Moreover, when these conditions are satisfied, the asphericity
structure $M_{as}$ in (i) is unique.

Of course, we'll say it is the asphericity structure on M \uline{genera-%
ted} by the full subcategory N of M, and the latter will be called a
\uline{generating} \uline{subcategory} for the given asphericity structure.

The proposition is a tautology, in view of the definitions and of
prop.1 (p. 232) and it's corollary. The f\unsure{orm} (ii) or (iii) of the condi-%
tion shows that it depends only upon the category N, not upon the way
in which N is embedded in a larger category M. We may call a category N
satisfying condition (iii) above an \uline{aspherator} , by which we would like
to express that this category represents a standart way of generating
asphericity structures, through \deleted{the}\inserted{any full} embedding of N into a category M.
This condition is automatically satisfied if N is small, more generally
it holds if N is equivalent to a small category. It should be kept in
mind that the condition depends both on the choice of \unsure{\over{th}{th}}e basic universe
$\underline{U}$ and on the choice of the basic localizer $\underline{W}$ in the corresponding large
category (Cat). It looks pretty sure that the condition is not always
satisfied (I doubt it is for N = (Sets) say), but I confess I did'nt
sit down to make an explicit example.

\margin{60}
